\clearpage
\VerdierPhysicalPageAnchor{0126}%
\begingroup
\small
\setlength{\emergencystretch}{1.5em}%
\setlength{\abovedisplayskip}{4pt plus 2pt minus 2pt}%
\setlength{\belowdisplayskip}{4pt plus 2pt minus 2pt}%

\VerdierSemanticAnchor{ch2-1-3-7}
\noindent\textbf{1.3.7.}\quad
\VerdierSemanticAnchor{ch2-1-3-7-multi-additive-functor}
Let
$F\colon \prod_{i\in[n]} C_i\longrightarrow C$
be a multi-additive functor. Suppose that countable products are
representable in $C$. It follows immediately from
(Chapter I, 3.4.3 and 3.4.4) that the $\Pi$-extension of the functor $F$
to complexes up to homotopy,

\VerdierSemanticAnchor{ch2-1-3-7-pi-extension}
\[
  K_{\Pi}F\colon \prod_{i\in[n]} K(C_i)\longrightarrow K(C),
\]

\VerdierSemanticAnchor{ch2-1-3-7-multi-exact}
\noindent is a multi-exact functor.

\VerdierSemanticAnchor{ch2-1-3-7-koszul-twist}
This functor is a $\mathbb Z^n$-functor twisted by the Koszul cocycle in
accordance with the sign rule (Chapter I, 1.6.6).

\VerdierSemanticAnchor{ch2-1-3-7-sigma-extension}
There are analogous results for the functor $K_{\Sigma}F$.

\VerdierSemanticAnchor{ch2-1-3-8}
\medskip
\noindent\textbf{1.3.8.}\quad
\VerdierSemanticAnchor{ch2-1-3-8-canonical-opposite-isomorphism}
Let $C$ be an additive category, and let
$\iota\colon K(C)^\circ\longrightarrow K(C^\circ)$ be the canonical
isomorphism $\iota=\iota_1$ (Chapter I, 2.5.13).

\VerdierSemanticAnchor{ch2-1-3-8-triangulated-opposite}
Equip the category $K(C)^\circ$ with the triangulated structure described
in (1.1.7).

\VerdierSemanticAnchor{ch2-1-3-8-exactness}
It follows from (Chapter I, 3.4.1) that $\iota$ is an exact functor
between triangulated categories. Consequently, $\iota$ is an isomorphism
of triangulated categories.

\VerdierSemanticAnchor{ch2-1-3-8-identification}
Henceforth, unless otherwise stated, we identify the triangulated
categories $K(C)^\circ$ and $K(C^\circ)$ by means of the isomorphism
$\iota$.
\endgroup
